\IEEEPARstart{E}{motion} plays a crucial role in human life, influencing decision-making processes and interpersonal interactions~\cite{COWIE20035}. With the rapid advancement of technology and the increasing integration of artificial intelligence into everyday tasks, there is a growing need to develop AI systems capable of understanding and responding to human emotions effectively~\cite{Venkataramanan2019}. In particular, intelligent assistants and conversational agents are expected to recognize emotional tone and generate responses that align with users’ affective states. However, relying solely on textual input is often insufficient, as human emotions are subtle, overlapping, and highly context-dependent, making them difficult to accurately capture using unimodal signals~\cite{PAYALA2026}. Furthermore, external factors such as background noise, variations in speaker tone, and environmental conditions introduce additional challenges for emotion recognition systems~\cite{PAYALA2026}. As a result, multimodal approaches that combine text and speech have gained attention due to their ability to capture complementary emotional cues from both linguistic and paralinguistic information. While visual signals, such as facial expressions, can further enhance emotion recognition performance, their practical adoption remains limited due to challenges in large-scale data collection, annotation, and real-world deployment constraints, including privacy concerns and user preferences~\cite{Lian2023}. In many real-world applications, users are more inclined to provide audio rather than video input, making audio–text multimodal systems a more feasible and scalable solution.

Emotion recognition has attracted increasing research attention in recent years, leading to significant advancements in both unimodal and multimodal learning frameworks. Early studies primarily focused on unimodal emotion recognition using either textual information\cite{Cortiz2022, Batbaatar2019, Han2023} or speech signals through handcrafted acoustic features\cite{Mustaqeem2020, Venkataramanan2019, Novais2023} and, more recently, self-supervised learned speech representations\cite{Atmaja2022, Yingzhi2021}. Although these approaches achieved promising performance, relying on a single modality often limits the ability of the model to capture the complete emotional context. To overcome these limitations, recent research has increasingly shifted toward multimodal emotion recognition by integrating complementary information from speech and text modalities. Existing multimodal approaches employ various fusion strategies, including feature interaction methods\cite{Koudounas2026, Chen2020, Chudasama2022, Dikbiyik2025, ISLAM2026, Kakuba2023}, cross-modal attention mechanisms\cite{Krishna2020}, knowledge distillation frameworks\cite{KIM2026}, and ensemble-based approaches that combine independently learned predictions or representations from different modalities\cite{PAYALA2026, Sebastian2019}. More recently, contrastive learning has emerged as an effective paradigm for multimodal representation learning. Existing contrastive multimodal approaches can generally be categorized into cross-modal alignment methods\cite{qian2022}, fusion-integrated contrastive learning frameworks \cite{Shi2024, Li2025, Kim2023}, and pretraining-based contrastive approaches\cite{Jing2026, Nguyen2024}.

Several studies have demonstrated the effectiveness of multimodal learning and contrastive objectives for emotion recognition. For example, the ablation study presented in\cite{Shi2024} showed that multimodal fusion significantly outperformed unimodal approaches on both the IEMOCAP and MELD datasets, achieving approximately 18\% improvement compared to speech-only models and around 4\% improvement compared to text-only models. Similarly, \cite{qian2022} reported that incorporating cross-modal contrastive alignment improved weighted accuracy on the IEMOCAP dataset by 1.44\%. In addition,\cite{Li2025} demonstrated that integrating label-aware contrastive learning within multimodal fusion increased the weighted average F1-score on the MELD dataset by approximately 4\%. These findings collectively highlight the effectiveness of combining multimodal learning with contrastive representation objectives for emotion recognition tasks.

Despite these advances, several challenges remain insufficiently addressed. Existing methods still struggle to achieve robust generalization across different datasets, speakers, and linguistic distributions, while many fusion strategies rely on static interactions that fail to adaptively model the varying reliability and importance of modalities across conversational contexts. Furthermore, current contrastive learning approaches often focus on either cross-modal alignment or discriminative representation learning independently, limiting their ability to jointly optimize semantic consistency and emotional separability within multimodal representations. To address these limitations, we propose a unified multimodal framework that integrates speech and text through a dynamic utterance-level fusion mechanism combined with structured cross-modal contrastive learning.

Specifically, the proposed framework introduces a chunk-based audio modeling strategy in which each utterance is segmented into overlapping temporal chunks to better capture fine-grained variations in emotional expression. Supervised contrastive learning is applied at the chunk level to enhance intra-class compactness and inter-class separability within the audio representation space. To further improve robustness, a confidence-aware aggregation mechanism is employed to weight chunk-level representations according to their reliability, producing a more stable utterance-level embedding. In parallel, a cross-modal contrastive objective aligns audio and text representations within a shared embedding space, improving inter-modal consistency. By jointly optimizing dynamic fusion, chunk-level supervised contrastive learning, cross-modal alignment, and confidence-aware aggregation, the proposed framework achieves improved generalization and provides a robust and discriminative solution for multimodal emotion recognition across diverse conditions.

\subsection*{Main Contributions:}

\begin{itemize}
\item We propose \textsc{Daf-Cl}, a unified multimodal emotion recognition framework that jointly learns dynamic audio–text fusion and cross-modal representation alignment. The proposed approach adaptively balances modality contributions at the utterance level, enabling robust performance under varying modality reliability and noise conditions.
\item We introduce a chunk-based audio representation learning strategy combined with supervised contrastive learning to capture fine-grained temporal variations in speech while enhancing intra-class compactness and inter-class separability.
\item We design a confidence-aware aggregation mechanism that adaptively weights chunk-level representations according to their reliability, producing more stable and robust utterance-level representations.
\item We formulate a unified multi-objective optimization framework that jointly integrates classification loss, cross-modal contrastive alignment, and supervised contrastive learning to enable end-to-end learning of discriminative and aligned multimodal representations.
\item We conduct extensive experiments and ablation studies across multiple datasets and cross-lingual settings, demonstrating the robustness and generalization capability of the proposed framework under diverse conditions.
\end{itemize}

\subsection*{Paper Organization:} The rest of the paper is organized as follows. In Section~\ref{sec:related work}, we review previous work related to our paper. In Section. In Section~\ref{sec:Model 1}, we develop our framework, \textsc{Daf-Cl}. Then, in Section~\ref{sec:performance evaluation}, . Finally, the main conclusions of this paper are summarized in Section~\ref{sec:Conclusion}.


